%% kapa: exam=1 round=31 year=2020 period=1 form=A subject=economics subject_ko=경제학원론 seq=25 number=65
%% source: Q-Net 감정평가사 기출문제 (cst00309) artlSeq=5207219 file=extracted/1008/5207219/A형/A형/2020년도 제31회 감정평가사 1차 시험_1교시_A형.pdf
%% conversion: kordoc 4.24.51 -> tools/gen1.py (auto); review=proofread-sonnet; flags=none
\begin{question}{65}
A국에서 인플레이션 갭과 산출량 갭이 모두 확대될 때, 테일러 준칙(Taylor's rule)에 따른 중앙은행의 정책은?
\begin{choices}
  \item 정책금리를 인상한다.
  \item 정책금리를 인하한다.
  \item 정책금리를 조정하지 않는다.
  \item 지급준비율을 인하한다.
  \item 지급준비율을 변경하지 않는다.
\end{choices}
\end{question}
\answer{1}
